\documentclass[a4paper,11pt]{article}
\usepackage[a4paper, top=2.5cm, bottom=2.5cm, left=2cm, right=2cm]{geometry}
\usepackage{enumitem}
\setlength\parindent{0pt}
\usepackage{fancyhdr}

\pagestyle{fancy}
\fancyhf{}
\fancyfoot[L]{TN Enabling Technology}
\fancyfoot[C]{\today}
\fancyfoot[R]{2025-2026}

\renewcommand{\footrulewidth}{0.4pt}
\renewcommand{\headrulewidth}{0pt}

\begin{document}

\section*{Opdrachtbeschrijving Memo 2 -- Technische Stage}
\subsection*{Voortgang en resultaten \textbf{(20\%)}}
Beschrijf de voortgang binnen jouw stage en de resultaten van uitgevoerde experimenten of analyses. Behandel minimaal:
\begin{itemize}[leftmargin=1.5cm, partopsep=1pt, itemsep=2pt, parsep=0pt]
    \item Welke voortgang is geboekt ten opzichte van vorige rapportages;
    \item Resultaten van experimenten of simulaties, helder gepresenteerd;
    \item Analyses onderbouwd en correct uitgevoerd;
    \item Toelichting op resultaten logisch en gestructureerd.
\end{itemize}

\textbf{Criterium:} De lezer moet een volledig en logisch beeld krijgen van de voortgang, resultaten en uitgevoerde analyses.

\subsection*{Discussie en volgende stappen \textbf{(20\%)}}
Bespreek de resultaten in relatie tot het stagedoel en beschrijf vervolgstappen. Behandel minimaal:
\begin{itemize}[leftmargin=1.5cm, partopsep=1pt, itemsep=2pt, parsep=0pt]
    \item Verklaring en interpretatie van de resultaten;
    \item Verband met het stagedoel en relevantie voor het project;
    \item Belangrijke inzichten en beperkingen;
    \item Conclusies onderbouwd en logisch;
    \item Concrete en haalbare vervolgstappen.
\end{itemize}

\textbf{Criterium:} De lezer moet kunnen volgen hoe resultaten bijdragen aan het stagedoel, welke inzichten zijn opgedaan en welke stappen nog gezet moeten worden.

\subsection*{Schematische tekening en toelichting \textbf{(20\%)}}
Maak een duidelijke schematische weergave van jouw experiment of opstelling. Behandel minimaal:
\begin{itemize}[leftmargin=1.5cm, partopsep=1pt, itemsep=2pt, parsep=0pt]
    \item Digitale, overzichtelijke en professionele schets van de opstelling;
    \item Representatieve weergave van de experimentele situatie;
    \item Volledige en begrijpelijke toelichting en onderschrift;
    \item Eventuele belangrijke opmerkingen over uitvoering of beperkingen.
\end{itemize}

\textbf{Criterium:} De schets en toelichting moeten het experiment volledig en correct weergeven, zodat een medestudent het kan begrijpen en reproduceren.

\subsection*{Opmaak figuren \textbf{(20\%)}}
Zorg dat figuren duidelijk en overzichtelijk zijn en ondersteunend voor de resultaten. Behandel minimaal:
\begin{itemize}[leftmargin=1.5cm, partopsep=1pt, itemsep=2pt, parsep=0pt]
    \item Professioneel opgemaakte figuren met correcte nummering, labels en onderschriften;
    \item Relevante en beknopte figuren, zonder overbodige data;
    \item Figuren en onderschriften ondersteunen de interpretatie van resultaten volledig.
\end{itemize}

\textbf{Criterium:} De figuren moeten overzichtelijk zijn, goed interpreteerbaar en correct gekoppeld aan de resultaten.

\subsection*{Algemene criteria \textbf{(20\%)}}
\begin{itemize}[leftmargin=1.5cm, partopsep=1pt, itemsep=2pt, parsep=0pt]
    \item Lengte: maximaal 4 A4;
    \item Grammaticaal en taalkundig correct;
    \item Nette en overzichtelijke lay-out;
    \item Op tijd ingeleverd;
    \item Geen plagiaat;
    \item Correcte verwijzingen en gebruik van AI (inclusief prompts).
\end{itemize}

\textbf{Criterium:} Het verslag voldoet aan algemene academische standaarden en is volledig en correct gepresenteerd.

\end{document}
